% Source PDF: source/普通高考/2012/2012北京文.pdf, page 1, question 4.
% Original flowchart is vector/text artwork; preserved PNG remains under img/.
% Nodes: 开始; k=0,S=1; k<3?; S=S·2^k; k=k+1; 输出S; 结束.
% Directed edges: 开始→初始化→判定; 判定(是)→S更新→k更新→共享入口→判定;
% 判定(否)→输出S→结束. The loop return and main downward arrow are separate.

\begin{tikzpicture}[
  >=Stealth,
  x=1cm,
  y=1cm,
  line width=0.45pt,
  line cap=round,
  line join=round,
  font=\normalsize,
  flowtext/.style={anchor=center, align=center,
    execute at begin node={\setlength{\parindent}{0pt}}},
  flowbox/.style={flowtext, draw, minimum height=0.68cm,
    inner xsep=4pt, inner ysep=2pt},
  terminal/.style={flowbox, rounded corners=2pt, minimum width=1.35cm},
  process/.style={flowbox, minimum width=3.20cm},
  decision/.style={flowbox, diamond, aspect=2.8, inner sep=0pt,
    minimum width=4.55cm, minimum height=1.08cm},
  io/.style={flowbox, trapezium, trapezium left angle=72,
    trapezium right angle=108, minimum height=0.78cm},
  branchlabel/.style={flowtext, inner sep=0pt},
]
  \node[terminal] (start) at (0,11.50) {开始};
  \node[process, minimum width=3.95cm] (init) at (0,9.80) {$k=0,S=1$};
  \node[decision] (test) at (0,4.60) {$k<3$};
  \node[process, minimum width=3.20cm] (smul) at (4.25,6.15)
    {$S=S\cdot 2^k$};
  \node[process, minimum width=3.05cm] (inc) at (4.25,8.05)
    {$k=k+1$};
  \node[io, minimum width=2.95cm] (output) at (0,2.55) {输出 $S$};
  \node[terminal] (finish) at (0,0.80) {结束};

  \coordinate (loopjoin) at (0,8.85);
  % Keep the return segment clear of init; its arrow ends at the loop-entry junction.
  \coordinate (loop-tip) at (0.45,8.85);
  \coordinate (loop-right) at (4.25,8.85);
  \coordinate (loop-bottom) at (4.25,4.60);

  \draw[->] (start.south) -- (init.north);
  \draw (init.south) -- (loopjoin);
  \draw[->] (loopjoin) -- (test.north);

  % 是 branch: update S, then update k, then return left into the shared entry.
  \draw (test.east) -- (loop-bottom);
  \draw[->] (loop-bottom) -- (smul.south);
  \draw[->] (smul.north) -- (inc.south);
  \draw (inc.north) -- (loop-right);
  \draw[->] (loop-right) -- (loop-tip) -- (loopjoin);
  \node[branchlabel] at (1.85,4.00) {是};

  % 否 branch remains an independent output path.
  \draw[->] (test.south) -- (output.north);
  \node[branchlabel] at (0.32,3.85) {否};
  \draw[->] (output.south) -- (finish.north);

  \path[use as bounding box] (-2.55,0.25) rectangle (5.05,11.90);
\end{tikzpicture}
